\chapter{Metodologi}
Tugas akhir ini menggunakan tiga pilar metodologis utama yaitu analisis ekonofisika, teori permainan strategis, dan komputasi kuantum variasional diintegrasikan untuk mengoptimalkan pemilihan portofolio biner secara hibrida. Untuk menghindari kebingungan simbolis, diberikan notasi matematis yang digunakan pada modul-modul metodologi sebagaimana pada Tabel \ref{tab:simbol_metodologi}.

\begin{table}[htbp]
\centering
\caption{Daftar Simbol dan Notasi Matematis dalam Metodologi.}
\label{tab:simbol_metodologi}
\begin{tabular}{|c|l|}
\hline
\textbf{Simbol} & \textbf{Definisi} \\ \hline
$\mu$ & Vektor ekspektasi imbal hasil (\textit{expected return}) aset. \\ \hline
$\Sigma$ & Matriks kovariansi yang merepresentasikan risiko sistemik antar aset. \\ \hline
$\gamma$ & Parameter penghindaran risiko (\textit{risk aversion}) endogen. \\ \hline
$\lambda$ & Koefisien pinalti untuk penegakan kendala kardinalitas aset. \\ \hline
$K$ & Jumlah target aset yang akan dipilih ke dalam portofolio (\textit{budget}). \\ \hline
$N$ & Jumlah total aset yang tersedia dalam semesta pemilihan. \\ \hline
$x_i$ & Variabel keputusan biner untuk aset ke-$i$ ($x_i \in \{0, 1\}$). \\ \hline
$\Phi$ & Fungsi potensial global dalam kerangka \textit{Exact Potential Game}. \\ \hline
$\hat{\mathcal{H}}$ & Operator Hamiltonian yang merepresentasikan energi total sistem kuantum. \\ \hline
$\theta$ & Vektor parameter sudut rotasi pada sirkuit kuantum variasional. \\ \hline
$a_k, c_k$ & Hiperparameter laju pembelajaran dan perturbasi pada optimasi SPSA. \\ \hline
\text{s.t.} & \textit{subject to} (dengan kendala atau batasan matematis). \\ \hline
\end{tabular}
\end{table}

\section{Alur Penelitian}
Tiga pilar metodologis utama yang mencakup analisis ekonofisika, teori permainan strategis, dan komputasi kuantum variasional diintegrasikan untuk mengoptimalkan pemilihan portofolio biner secara hibrida. Alur penelitian dimulai dengan mengekstraksi data historis saham indeks LQ45 yang kemudian ditransformasi menjadi variabel imbal hasil logaritmik guna mendeteksi korelasi stokastik antar aset finansial. Interaksi strategis antar aset tersebut dipetakan menggunakan parameter \textit{Exact Potential Game} (EPG) untuk mengidentifikasi struktur kestabilan Nash yang berfungsi sebagai input utama bagi Hamiltonian Ising. Proses pencarian solusi optimal dieksekusi melalui algoritma \textit{Hardware-Efficient Variational Quantum Eigensolver} (HE-VQE) dengan memanfaatkan optimasi SPSA untuk meminimalkan fungsi energi sistem pada perangkat keras kuantum. Rangkaian penelitian ditutup dengan melakukan simulasi \textit{backtesting} secara komprehensif untuk mengevaluasi performa model terhadap data pasar riil melalui metrik \textit{Sharpe Ratio} dan \textit{Maximum Drawdown}.

Struktur operasional penelitian dirancang ke dalam empat modul sekuensial yang saling berinteraksi secara dinamis guna menjamin integritas data dan akurasi solusi optimasi yang dihasilkan. Modul pertama difokuskan pada akuisisi serta pra-pemrosesan data mentah dan modul kedua diarahkan untuk mengekstraksi parameter interaksi melalui pendekatan teori informasi kuantum serta ekonofisika sistemik. Modul ketiga diposisikan sebagai inti optimasi yang melibatkan konstruksi matriks QUBO dan pelatihan parameter sirkuit kuantum menggunakan metode \textit{Adaptive Depth} untuk mencapai konvergensi energi minimum. Modul keempat ditugaskan untuk mengeksekusi strategi investasi secara virtual serta menghitung metrik performa portofolio guna memberikan gambaran objektif mengenai profitabilitas dan tingkat risiko model. Seluruh tahapan kerja ini disinkronkan untuk memastikan bahwa setiap transisi data memiliki landasan teoretis yang kuat berdasarkan prinsip mekanika kuantum dan dinamika pasar modal kontemporer.



\begin{figure}[ht]
    \centering
    \begin{tikzpicture}[node distance=1.3cm, every node/.style={fill=white, font=\footnotesize}, align=center]
        % Styles
        \tikzstyle{startstop} = [ellipse, minimum width=2.5cm, minimum height=0.6cm, text centered, draw=black, fill=white]
        \tikzstyle{process} = [rectangle, minimum width=3.2cm, minimum height=0.8cm, text centered, draw=black, fill=white]
        \tikzstyle{decision} = [diamond, minimum width=2.5cm, minimum height=1cm, text centered, draw=black, fill=white, aspect=2]
        \tikzstyle{connector} = [circle, draw=black, minimum size=0.5cm, text centered, inner sep=2pt]
        \tikzstyle{arrow} = [thick,->,>=stealth]
        
        % KOLOM 1: Pra-pemrosesan & Parameterisasi
        \node (start) [startstop] {Mulai};
        \node (data) [process, below of=start, yshift=-0.5cm] {Akuisisi Data \& \\ Setup Jendela Waktu};
        \node (param) [process, below of=data, yshift=-0.5cm] {Hitung $\mu$, $\Sigma$ \& \\ \textit{Risk Aversion} endogen ($\gamma$)};
        \node (hamil) [process, below of=param, yshift=-0.5cm] {Formulasi QUBO \& \\ Hamiltonian ($h_i, J_{ij}$)};
        \node (nash) [process, below of=hamil, yshift=-0.5cm] {Pencarian \\ \textit{Nash Equilibrium (SBR)}};
        \node (c1a) [connector, below of=nash, yshift=-0.5cm] {1};
        \node (c3b) [connector, left of=data, xshift=-1cm] {3};
        
        % KOLOM 2: Optimasi Kuantum (VQE)
        \node (c1b) [connector, right of=start, xshift=4.0cm] {1};
        \node (vqe_init) [process, below of=c1b, yshift=-0.2cm] {Inisialisasi State Awal \\ $|0\dots0\rangle \to |\text{bitstring}_{NE}\rangle$ \\ \& Parameter $\theta$};
        \node (spsa) [process, below of=vqe_init, yshift=-0.5cm] {Optimasi SPSA};
        \node (conv) [decision, below of=spsa, yshift=-0.5cm] {Konvergen?};
        \node (layer_eval) [process, below of=conv, yshift=-0.6cm] {Simpan Energi $E_d$ \\ pada \textit{Depth} $d$};
        \node (layer_dec) [decision, below of=layer_eval, yshift=-0.5cm] {$d < D_{max}$?};
        \node (best_state) [process, below of=layer_dec, yshift=-0.7cm] {Pilih Bitstring \\ dengan $E_{min}$};
        \node (c2a) [connector, below of=best_state, yshift=-0.3cm] {2};
        
        % KOLOM 3: Evaluasi & Backtesting
        \node (c2b) [connector, right of=c1b, xshift=4.0cm] {2};
        \node (rebal) [process, below of=c2b, yshift=-0.5cm] {Eksekusi \textit{Rebalancing} \\ Portofolio};
        \node (end_dec) [decision, below of=rebal, yshift=-1cm] {Akhir \\ Periode?};
        \node (eval) [process, below of=end_dec, yshift=-1cm] {Kalkulasi Metrik \\ Kinerja Akhir};
        \node (stop) [startstop, below of=eval, yshift=-0.5cm] {Selesai};
        \node (c3a) [connector, right of=end_dec, xshift=2.1cm] {3};

        % Alur Kolom 1
        \draw [arrow] (start) -- (data);
        \draw [arrow] (data) -- (param);
        \draw [arrow] (param) -- (hamil);
        \draw [arrow] (hamil) -- (nash);
        \draw [arrow] (nash) -- (c1a);
        \draw [arrow] (c3b) -- (data);

        % Alur Kolom 2
        \draw [arrow] (c1b) -- (vqe_init);
        \draw [arrow] (vqe_init) -- (spsa);
        \draw [arrow] (spsa) -- (conv);
        \draw [arrow] (conv) -- node[anchor=east] {Ya} (layer_eval);
        \draw [arrow] (conv.west) -- ++(-0.6,0) |- (spsa.west) node[pos=0.25, anchor=west] {Tidak};
        \draw [arrow] (layer_eval) -- (layer_dec);
        \draw [arrow] (layer_dec) -- node[anchor=east] {Tidak} (best_state);
        \draw [arrow] (layer_dec.east) -- ++(0.7,0) |- (vqe_init.east) node[pos=0.25, anchor=west] {Ya (Tambah \\ \textit{Depth})};
        \draw [arrow] (best_state) -- (c2a);

        % Alur Kolom 3
        \draw [arrow] (c2b) -- (rebal);
        \draw [arrow] (rebal) -- (end_dec);
        \draw [arrow] (end_dec) -- node[anchor=east] {Ya} (eval);
        \draw [arrow] (eval) -- (stop);
        \draw [arrow] (end_dec.east) -- (c3a.west) node[midway, anchor=south] {Tidak};

    \end{tikzpicture}
    \caption{Diagram Alir Metodologi Optimasi Portofolio Berbasis GT-VQE (Nash SBR dan Adaptive VQE dengan Penalty Annealing).}
    \label{fig:flowchart_gt_vqe}
\end{figure}

\section{Akuisisi dan Pra-pemrosesan Data}

Data primer penelitian ini diperoleh melalui pengunduhan harga penutupan harian (\textit{closing price}) dari seluruh emiten yang terdaftar dalam indeks saham IHSG menggunakan modul akuisisi data yang terhubung langsung dengan API \textit{Yahoo Finance}. Rentang waktu pengambilan data ditetapkan mencakup periode observasi yang cukup panjang guna memastikan bahwa model mampu menangkap berbagai dinamika pasar mulai dari fase pertumbuhan hingga fase koreksi sistemik yang ekstrem. Data mentah yang telah berhasil diperoleh disimpan ke dalam struktur basis data sementara untuk mempermudah proses sanitasi serta pembersihan dari nilai kosong atau anomali data akibat aksi korporasi yang tidak konsisten. Indeks saham IHSG dipilih sebagai objek penelitian utama karena memiliki tingkat likuiditas yang sangat tinggi serta kapitalisasi pasar yang besar sehingga hasil optimasi memiliki relevansi praktis yang kuat bagi para investor institusional di Indonesia.

Harga penutupan absolut ditransformasi menjadi variabel imbal hasil logaritmik (\textit{log returns}) menggunakan modul perhitungan imbal hasil guna mencapai kondisi stasioneritas data secara statistik yang sangat diperlukan dalam pemodelan deret waktu finansial. Penggunaan imbal hasil logaritmik diprioritaskan dibandingkan dengan imbal hasil sederhana karena memiliki sifat aditivitas yang memungkinkan dilakukannya agregasi temporal secara linear tanpa mengalami distorsi nilai yang signifikan pada proyeksi jangka panjang. Ekspektasi imbal hasil harian $\mu$ dan matriks kovariansi $\Sigma$ dihitung dari data \textit{log return} tersebut sebagai basis input fundamental yang akan dipetakan ke dalam parameter interaksi fisis pada model Hamiltonian. Perhitungan \textit{volatility drag} diintegrasikan secara eksplisit untuk mengonversi imbal hasil aritmatika menjadi imbal hasil logaritmik yang lebih akurat dalam merepresentasikan laju pertumbuhan modal secara geometris di tengah fluktuasi pasar.

Derajat penghindaran risiko (\textit{risk aversion}) $\gamma$ ditentukan secara endogen menggunakan modul parameterisasi risiko yang secara otomatis menyesuaikan nilai parameter tersebut berdasarkan kondisi volatilitas pasar pada setiap jendela waktu observasi. Nilai $\gamma$ dihitung melalui transformasi logistik terhadap skor-Z agregat yang merepresentasikan \textit{Sharpe Ratio} lintas aset dalam jendela waktu tersebut:
\begin{equation}
    Z = \frac{\bar{\mu}}{\bar{\sigma}}, \quad \gamma = \frac{1}{1 + e^{Z}}
\end{equation}
di mana $\bar{\mu}$ adalah rata-rata magnitudo imbal hasil dan $\bar{\sigma}$ adalah rata-rata varians aset. Pendekatan parameterisasi endogen ini memastikan bahwa model optimasi tidak bersifat kaku melainkan mampu merespons lonjakan risiko sistemik secara cepat melalui penyesuaian bobot penalti pada fungsi objektif portofolio yang sedang dibangun. Seluruh alur kerja akuisisi data, transformasi statistik, hingga ekstraksi parameter $\gamma$ dirangkum ke dalam sebuah prosedur sistematis yang disajikan secara formal pada Algoritma \ref{alg:data_prep}.


\begin{algorithm}[ht]
\caption{Akuisisi dan Pra-pemrosesan Data Finansial}
\label{alg:data_prep}
\begin{algorithmic}[1]
\Procedure{DataPreprocessing}{tickers, $t_{start}, t_{end}, K$}
    \State $\mathbf{P} \gets \text{DownloadClosingPrice}(tickers, t_{start}, t_{end})$
    \State $\mathbf{P}_{clean} \gets \text{DropNaN}(\mathbf{P})$
    \State $r_{log} \gets \ln(\mathbf{P}_{t} / \mathbf{P}_{t-1})$
    \State $\bar{\mu} \gets \text{Mean}(|r_{log}|)$
    \State $\bar{\sigma} \gets \text{Mean}(\text{Variance}(r_{log}))$
    \State $Z \gets \bar{\mu} / \bar{\sigma}$ \Comment{Skor-Z Agregat}
    \State $\gamma \gets 1 / (1 + \exp(Z))$ \Comment{Risk Aversion Endogen}
    \State $\mu_{log}, \Sigma \gets \text{CalculateStats}(r_{log})$
    \State \Return $\mu_{log}, \Sigma, \gamma$
\EndProcedure
\end{algorithmic}
\end{algorithm}

\section{Pemodelan Hamiltonian Portofolio dan Pemetaan Qubit}

Masalah optimasi portofolio biner dirumuskan ke dalam format \textit{Quadratic Unconstrained Binary Optimization} (QUBO) yang mengintegrasikan fungsi biaya objektif dan fungsi penalti kendala ke dalam satu ekspresi matematis tunggal. Komponen objektif dirancang untuk meminimalkan risiko portofolio yang direpresentasikan oleh matriks kovariansi $\Sigma$ sekaligus memaksimalkan ekspektasi imbal hasil $\mu$ yang telah disesuaikan dengan parameter penghindaran risiko $\gamma$. Komponen penalti diterapkan secara ketat guna menegakkan batasan kardinalitas aset $K$ sehingga sistem hanya akan memberikan solusi valid yang memilih tepat sejumlah aset yang diinginkan dari total $N$ aset yang tersedia. Seluruh koefisien dalam matriks kuadratik $Q$ disusun secara sistematis menggunakan logika yang diimplementasikan pada modul strategi investasi guna menjamin bahwa lanskap energi sistem mencerminkan utilitas finansial yang akurat.

Formulasi QUBO ditransformasikan menjadi model Ising melalui pemetaan variabel biner $x_i \in \{0, 1\}$ ke dalam variabel \textit{spin} fisis $z_i \in \{+1, -1\}$ menggunakan substitusi linear $x_i = (1 - z_i)/2$ pada setiap variabel keputusan. Transformasi ini dilakukan untuk menghasilkan parameter bias individu $h_i$ dan parameter interaksi antar aset $J_{ij}$ yang secara kolektif membentuk lanskap energi total dari sistem Hamiltonian yang akan dioptimalkan. Proses konversi ini dieksekusi secara terpisah untuk komponen objektif dan komponen penalti guna memudahkan proses kalibrasi bobot pinalti $\lambda$ terhadap skala energi imbal hasil dan risiko yang dinamis. Modul konstruksi Hamiltonian digunakan untuk menyusun operator Hamiltonian tersebut ke dalam jumlahan produk operator \textit{Pauli-Z} yang dapat dievaluasi secara langsung oleh sirkuit kuantum pada basis komputasi standar.

Setiap aset tunggal yang terdaftar dalam indeks IHSG dipetakan secara satu-ke-satu ke dalam sebuah qubit di mana status dasar $|1\rangle$ memberikan interpretasi fisis bahwa aset tersebut terpilih masuk ke dalam keranjang portofolio biner. Nilai ekspektasi energi yang terukur dari sirkuit kuantum dipastikan memiliki korespondensi linear terhadap nilai risiko dan imbal hasil riil dari konfigurasi aset yang sedang dianalisis oleh algoritma. Seluruh arsitektur pemetaan ini diintegrasikan agar model optimasi kuantum siap memasuki fase pencarian solusi melalui algoritma variasional yang memanfaatkan landasan teoretis dari mekanika kuantum statistik. Rincian operasional dari pembentukan matriks QUBO hingga menjadi operator Hamiltonian kuantum yang utuh disajikan pada Algoritma \ref{alg:hamiltonian_construction}.

\begin{algorithm}[ht]
\caption{Konstruksi Hamiltonian Portofolio Ising}
\label{alg:hamiltonian_construction}
\begin{algorithmic}[1]
\Procedure{BuildHamiltonian}{$\mu, \Sigma, \gamma, K, \lambda$}
    \State $Q_{ii} \gets \frac{\gamma \Sigma_{ii}}{2K^2} - \frac{\mu_i}{K} + \lambda(1 - 2K)$ \Comment{Diagonal QUBO}
    \State $Q_{ij} \gets \frac{\gamma \Sigma_{ij}}{2K^2} + \lambda$ \Comment{Off-diagonal QUBO}
    \State $h_i \gets -\frac{1}{2} (Q_{ii} + \sum_{j \neq i} Q_{ij})$ \Comment{Transformasi ke Bias Ising}
    \State $J_{ij} \gets \frac{1}{4} Q_{ij}$ \Comment{Transformasi ke Interaksi Ising}
    \State $C \gets \frac{1}{4} \sum_{i} Q_{ii} + \frac{1}{8} \sum_{i \neq j} Q_{ij} + \lambda K^2$ \Comment{Konstanta Energi}
    \State $H \gets \sum_{i} h_i Z_i + \sum_{i < j} J_{ij} Z_i Z_j + C \cdot I$ \Comment{Konstruksi Matriks Hamiltonian}
    \State \Return $H$
\EndProcedure
\end{algorithmic}
\end{algorithm}


\section{Algoritma Nash Sequential Best Response (SBR)}

Interaksi strategis antar aset finansial dimodelkan sebagai sebuah \textit{Exact Potential Game} (EPG) di mana perubahan utilitas individu dari setiap aset selaras dengan perubahan pada fungsi potensial global $\Phi$. Fungsi potensial ini dirancang sebagai representasi dari gabungan ekspektasi imbal hasil kolektif dan mitigasi risiko sistemik yang memiliki korespondensi matematis langsung terhadap nilai energi negatif pada Hamiltonian Ising. Dalam kerangka kerja EPG, setiap titik maksimum lokal dari fungsi potensial merupakan sebuah titik ekuilibrium Nash (\textit{Nash Equilibrium}) yang memberikan konfigurasi portofolio paling stabil secara strategis bagi seluruh aset yang terlibat. Landasan teoretis ini dimanfaatkan untuk menyederhanakan proses pencarian solusi stabil pada sistem multi-aset yang kompleks menjadi sebuah masalah optimasi fungsi potensial tunggal yang lebih efisien secara komputasi.

Algoritma \textit{Sequential Best Response} (SBR) diimplementasikan melalui modul pencarian ekuilibrium Nash untuk mengekstraksi titik ekuilibrium Nash tersebut secara iteratif melalui eksplorasi ruang solusi portofolio. Algoritma ini dirancang agar bekerja dengan melakukan operasi pertukaran (\textit{swap move}) antara satu aset yang telah terpilih di dalam portofolio dengan satu aset yang berada di luar portofolio yang mampu memberikan peningkatan nilai potensial paling signifikan. Proses iterasi dipastikan terus berlanjut hingga sistem mencapai kondisi \textit{steady-state} di mana tidak ditemukan lagi pertukaran aset tunggal yang dapat meningkatkan utilitas kolektif dari keranjang investasi. Solusi bitstring final yang diperoleh dari tahap pencarian Nash ini digunakan sebagai \textit{warm-start} untuk menginisialisasi sirkuit kuantum pada fase \textit{Variational Quantum Eigensolver} guna mempercepat proses konvergensi energi.

Mekanisme pencarian ini diterapkan secara konsisten pada setiap awal periode investasi guna memberikan titik awal pencarian yang memiliki makna finansial yang kuat bagi sistem kuantum. Jembatan integrasi antara teori permainan dan komputasi kuantum dirancang untuk memitigasi risiko terjebak dalam solusi suboptimal serta mengatasi fenomena dataran tandus yang sering menghambat efisiensi algoritma optimasi berbasis gradien. Setiap transisi variabel dari ranah strategi ekonomi ke ranah optimasi kuantum dipastikan memiliki landasan logis yang dapat dipertanggungjawabkan secara saintifik melalui pendekatan ekonofisika sistemik. Rincian prosedural dari algoritma pencarian ekuilibrium Nash menggunakan metode \textit{Sequential Best Response} tersebut disajikan secara rinci pada Algoritma \ref{alg:nash_sbr}.


\begin{algorithm}[ht]
\caption{Pencarian Nash Equilibrium menggunakan Sequential Best Response (SBR)}
\label{alg:nash_sbr}
\begin{algorithmic}[1]
\Procedure{NashSBR}{$\mu, \Sigma, \gamma, N, K, max\_iter$}
    \State $x \gets \text{InitialSelection}(K)$ \Comment{Inisialisasi $K$ aset pertama}
    \State $\Phi \gets \frac{1}{K} \sum \mu_i x_i - \frac{\gamma}{2K^2} \sum \Sigma_{ij} x_i x_j$ \Comment{Fungsi Potensial}
    \For{$iter = 1$ to $max\_iter$}
        \State $improved \gets \text{False}$
        \State $S_{in} \gets \{i \mid x_i = 1\}, S_{out} \gets \{j \mid x_j = 0\}$
        \For{$i \in S_{in}$ and $j \in S_{out}$}
            \State $x' \gets \text{Swap}(x, i, j)$
            \State $\Phi' \gets \text{CalculatePotential}(x', \mu, \Sigma, \gamma)$
            \If{$\Phi' > \Phi$}
                \State $x \gets x', \Phi \gets \Phi', improved \gets \text{True}$
                \State \textbf{break} \Comment{Pembaruan Sekuensial}
            \EndIf
        \EndFor
        \If{not $improved$} \textbf{break} \EndIf
    \EndFor
    \State \Return $x$ \Comment{Bitstring Ekuilibrium Nash}
\EndProcedure
\end{algorithmic}
\end{algorithm}


\section{Optimasi Variasional HE-VQE Adaptif}

Proses optimasi portofolio diimplementasikan melalui skema \textit{Hardware-Efficient Variational Quantum Eigensolver} (HE-VQE) hibrida yang mengombinasikan kekuatan komputasi kuantum dengan efisiensi pembaruan parameter klasik. Arsitektur sirkuit \textit{EfficientSU2} diadopsi sebagai ansatz utama yang terdiri dari lapisan gerbang rotasi qubit $R_y$ dan $R_z$ yang dapat diparameterisasi serta lapisan gerbang keterbelitan CNOT untuk memodelkan korelasi antar aset finansial. Strategi \textit{Adaptive Depth} diterapkan di dalam modul optimasi kuantum adaptif di mana kedalaman sirkuit ditingkatkan secara bertahap untuk mencari keseimbangan optimal antara ekspresivitas sirkuit dengan kerentanan terhadap akumulasi derau perangkat keras. Setiap penambahan lapisan baru dipastikan dilakukan melalui mekanisme \textit{warm-start} parameter agar energi sistem tetap stabil dan terhindar dari fenomena dataran tandus yang sering menghambat efisiensi pelatihan pada sistem kuantum.

Pembaruan parameter sirkuit dilakukan secara iteratif menggunakan optimizer \textit{Simultaneous Perturbation Stochastic Approximation} (SPSA) yang memiliki ketangguhan tinggi dalam menghadapi lanskap energi yang bersifat stokastik dan berderau. Strategi \textit{Penalty Annealing} diintegrasikan di mana koefisien pinalti $\lambda$ ditingkatkan secara moderat sepanjang proses iterasi guna memastikan bahwa sistem kuantum secara bertahap dipandu menuju ruang solusi yang memenuhi kendala kardinalitas aset. Inisialisasi awal parameter rotasi sirkuit dilakukan menggunakan bitstring yang diperoleh dari titik ekuilibrium Nash (\textit{NE-VQE}) untuk memberikan landasan kestabilan awal yang memiliki makna finansial yang kuat bagi sistem. Prosedur kalibrasi otomatis juga diterapkan untuk mencari \textit{learning rate} yang paling optimal menggunakan modul kalibrasi hyperparameter sebelum memulai fase optimasi utama guna menjamin kecepatan konvergensi yang maksimal.

Pengambilan sampel terhadap distribusi probabilitas status qubit (\textit{bitstrings}) dilakukan setelah proses optimasi mencapai titik konvergensi yang stabil untuk mengekstraksi konfigurasi portofolio terbaik yang dihasilkan oleh sistem. Pemilihan bitstring dengan nilai probabilitas tertinggi yang memenuhi batasan jumlah aset $K$ diprioritaskan sebagai solusi final yang diimplementasikan pada setiap jendela waktu investasi yang sedang dianalisis secara mendalam. Seluruh siklus optimasi hibrida ini dirancang agar bersifat robust terhadap fluktuasi data pasar sekaligus efisien dalam penggunaan sumber daya qubit yang terbatas pada perangkat keras kuantum era NISQ. Alur kerja sistematis dari algoritma VQE adaptif yang mencakup inisialisasi hangat, pinalti dinamis, hingga optimasi parameter menggunakan SPSA dirangkum secara formal pada Algoritma \ref{alg:vqe_adaptive}.

\begin{algorithm}[ht]
\caption{Adaptive HE-VQE dengan Penalty Annealing dan Nash Warm-Start}
\label{alg:vqe_adaptive}
\begin{algorithmic}[1]
\Procedure{AdaptiveVQE}{$H_{obj}, H_{pen}, x_{NE}, D_{max}, \lambda_{target}$}
    \State $\theta \gets \text{InitialParameters}(x_{NE})$ \Comment{Inisialisasi NE-VQE}
    \For{$d = 1$ to $D_{max}$}
        \State $a, c \gets \text{CalibrateSPSA}(H_{obj}, H_{pen}, \lambda_{target})$ \Comment{Kalibrasi \textit{Learning Rate}}
        \For{$k = 1$ to $N_{iter}$}
            \State $\lambda_k \gets \lambda_{target} \cdot (0.1 + 0.9 \cdot \frac{k}{0.7 N_{iter}})$ \Comment{Penalty Annealing}
            \State $\mathcal{L}(\theta) \gets \langle \psi(\theta) | H_{obj} + \lambda_k H_{pen} | \psi(\theta) \rangle$
            \State $g_k \gets \text{SPSA\_Gradient}(\mathcal{L}, \theta, c_k)$
            \State $\theta \gets \theta - a_k g_k$ \Comment{Update Parameter}
        \EndFor
        \State $\theta \gets \text{LayerExpansion}(\theta)$ \Comment{Tambah \textit{Depth} (Adaptif)}
        \If{$\text{Converged}(\theta)$} \textbf{break} \EndIf
    \EndFor
    \State $probs \gets |\langle x | \psi(\theta_{final}) \rangle|^2$ \Comment{Pengukuran Probabilitas}
    \State $x^* \gets \text{Argmax}(probs) \text{ s.t. } \sum x_i = K$
    \State \Return $x^*$
\EndProcedure
\end{algorithmic}
\end{algorithm}


\section{Simulasi Backtesting dan Metrik Evaluasi}

Performa model dievaluasi melalui prosedur \textit{sliding window backtesting} yang mensimulasikan keputusan investasi riil secara kronologis menggunakan modul simulasi investasi pada data historis saham LQ45. Setiap jendela waktu observasi dirancang agar terdiri dari fase pelatihan selama 126 hari bursa (\textit{lookback window}) untuk ekstraksi parameter finansial serta fase implementasi portofolio untuk periode investasi berikutnya yang telah ditentukan. Proses ini dijalankan secara berulang dan iteratif dengan menggeser jendela waktu ke depan secara periodik sehingga sirkuit kuantum selalu dilatih ulang menggunakan informasi pasar paling mutakhir guna menangkap pergeseran struktur korelasi antar aset. Variabel biaya transaksi diperhitungkan secara eksplisit pada setiap tahapan penyeimbangan portofolio (\textit{rebalancing}) guna memberikan estimasi imbal hasil bersih yang lebih realistis dan akurat sesuai dengan kondisi perdagangan sebenarnya di bursa saham.

Metrik kinerja utama yang meliputi \textit{Cumulative Return}, \textit{Sharpe Ratio}, serta \textit{Maximum Drawdown} dihitung melalui modul kalkulasi metrik performa untuk mendapatkan gambaran profil risiko dan imbal hasil model secara utuh. Metrik \textit{Sharpe Ratio} dimanfaatkan sebagai indikator efisiensi portofolio dalam menghasilkan imbal hasil tambahan untuk setiap unit risiko yang diambil oleh investor di tengah fluktuasi harga pasar yang bersifat stokastik. Metrik \textit{Maximum Drawdown} digunakan untuk memberikan wawasan mendalam mengenai potensi kerugian terbesar yang mungkin dialami dari titik puncak kekayaan selama periode investasi tertentu sebagai indikator ketahanan utama dari strategi yang diusulkan. Seluruh hasil perhitungan metrik ini dibandingkan secara langsung dengan performa indeks acuan (\textit{benchmark}) pasar untuk memvalidasi secara empiris keunggulan serta nilai tambah dari penggunaan teknologi komputasi kuantum dalam optimasi portofolio.

Seluruh modul metodologi mulai dari akuisisi data, pemodelan Hamiltonian, hingga optimasi kuantum diintegrasikan ke dalam sebuah rangkaian kerja global yang otomatis dan terstruktur secara sistematis. Setiap transisi data antar modul dipastikan memiliki validitas teoretis yang kuat dan didukung oleh hasil simulasi yang dapat direproduksi secara konsisten melalui repositori kode \texttt{GTQuantumInvest}. Seluruh rangkaian simulasi \textit{backtesting} yang mencakup siklus pembaruan portofolio hingga kalkulasi metrik kinerja akhir disusun ke dalam sebuah algoritma komprehensif yang disajikan pada Algoritma \ref{alg:backtesting_global}. Alur kerja ini dirancang untuk memberikan standar evaluasi yang objektif dan transparan bagi pengembangan algoritma finansial berbasis kuantum di masa depan agar memiliki daya saing yang tinggi pada lingkungan pasar modal yang dinamis.


\begin{algorithm}[ht]
\caption{Prosedur Backtesting Global Portofolio GT-VQE}
\label{alg:backtesting_global}
\begin{algorithmic}[1]
\Procedure{BacktestLoop}{$Data, K, \text{window}=126$}
    \State $\text{Timeline} \gets \text{GenerateInvestmentDates}(Data)$
    \For{$t \in \text{Timeline}$}
        \State $Data_{train} \gets Data[t-\text{window} : t]$
        \State $\mu, \Sigma, \gamma \gets \text{DataPreprocessing}(Data_{train}, K)$ \Comment{Algoritma \ref{alg:data_prep}}
        \State $H_{obj}, H_{pen} \gets \text{BuildHamiltonian}(\mu, \Sigma, \gamma, K, \lambda)$ \Comment{Algoritma \ref{alg:hamiltonian_construction}}
        \State $x_{NE} \gets \text{NashSBR}(\mu, \Sigma, \gamma, K)$ \Comment{Algoritma \ref{alg:nash_sbr}}
        \State $x^*_t, \theta_t \gets \text{AdaptiveVQE}(H_{obj}, H_{pen}, x_{NE}, D_{max})$ \Comment{Optimasi VQE Adaptif}
        \State $R_t \gets \text{ExecuteRebalance}(x^*_t, Data[t : t+1])$ \Comment{Eksekusi \textit{Rebalancing}}
        \State $\text{LogPortfolioHistory}(x^*_t, R_t, \theta_t)$
    \EndFor
    \State $\text{Metrics} \gets \text{ComputePerformance}(\text{History})$ \Comment{Kalkulasi Metrik Performa}
    \State \Return $\text{EquityCurve, Metrics}$
\EndProcedure
\end{algorithmic}
\end{algorithm}


\section{Optimasi Klasik SLSQP (\textit{Sequential Least Squares Programming})}

Sebagai standar pembanding (\textit{benchmark}) terhadap performa algoritma kuantum, penelitian ini menggunakan metode optimasi klasik \textit{Sequential Least Squares Programming} (SLSQP) untuk menyelesaikan masalah Markowitz kontinu. SLSQP merupakan algoritma optimasi iteratif yang dirancang untuk menangani masalah pemrograman non-linear dengan batasan kesamaan (\textit{equality constraints}) maupun batasan pertidaksamaan (\textit{inequality constraints}) secara efisien. Algoritma ini bekerja dengan merumuskan serangkaian sub-masalah \textit{Quadratic Programming} (QP) pada setiap langkah iterasi menggunakan aproksimasi deret Taylor orde kedua terhadap fungsi Lagrangian.

Proses optimasi dimulai dengan inisialisasi bobot portofolio awal yang seragam (\textit{uniform weights}) serta aproksimasi matriks Hessian sebagai matriks identitas. Pada setiap iterasi, gradien dari fungsi objektif dihitung dan matriks Hessian diperbarui menggunakan skema pembaruan BFGS (\textit{Broyden-Fletcher-Goldfarb-Shanno}) berdasarkan perubahan posisi dan gradien dari langkah sebelumnya. Arah pencarian optimal ditentukan dengan menyelesaikan sub-masalah QP yang mematuhi linierisasi dari batasan portofolio, yakni total bobot harus setara dengan satu dan setiap bobot berada dalam rentang $[0, 1]$. Algoritma ini dipastikan berkonvergensi menuju titik minimum lokal yang, untuk fungsi Markowitz yang bersifat konveks, juga merepresentasikan solusi optimal global. Rincian operasional dari algoritma SLSQP yang diimplementasikan disajikan pada Algoritma \ref{alg:slsqp_optimizer}.

\begin{algorithm}[ht]
\caption{Optimasi Klasik Markowitz menggunakan SLSQP}
\label{alg:slsqp_optimizer}
\begin{algorithmic}[1]
\Procedure{SLSQPOptimizer}{$\mu, \Sigma, \gamma$}
    \State $w_0 \gets \text{InitialWeights}(1/N)$ \Comment{Inisialisasi Seragam}
    \State $B_0 \gets \mathbf{I}$ \Comment{Inisialisasi Hessian Identitas}
    \For{$k = 0, 1, 2, \dots$ until converged}
        \State $g_k \gets \nabla f(w_k) = \gamma \Sigma w_k - \mu$ \Comment{Hitung Gradien}
        \State Solve Sub-QP: $\min_{d} (g_k^T d + \frac{1}{2} d^T B_k d) \text{ s.t. } \mathbf{1}^T d = 0, \text{ bounds}$
        \State $\alpha_k \gets \text{LineSearch}(w_k, d)$ \Comment{Cari ukuran langkah}
        \State $w_{k+1} \gets w_k + \alpha_k d$ \Comment{Update Bobot}
        \State $s_k \gets w_{k+1} - w_k, y_k \gets g_{k+1} - g_k$
        \State $B_{k+1} \gets \text{UpdateBFGS}(B_k, s_k, y_k)$ \Comment{Update Aproksimasi Hessian}
    \EndFor
    \State \Return $w_{final}$
\EndProcedure
\end{algorithmic}
\end{algorithm}


\section{Validasi Solusi Eksak Brute Force}
Metode pencarian solusi eksak berbasis \textit{brute force} diterapkan sebagai instrumen verifikasi absolut untuk menentukan kebenaran mutlak dari titik minimum global pada model energi Ising. Prosedur penelusuran solusi mutlak melalui perhitungan fungsi biaya Hamiltonian pada seluruh elemen ruang keadaan biner diformalkan ke dalam kerangka Algoritma \ref{alg:brute_force}. Algoritma komputasi eksak ini menerima parameter bias linier dan parameter kopling kuadratik untuk mengeksekusi perhitungan utilitas energi secara langsung pada setiap iterasi kombinasi aset. Mekanisme pencatatan nilai minimum dijalankan dengan memperbarui variabel energi terbaik setiap kali algoritma menemukan konfigurasi aset dengan ukuran energi yang lebih rendah secara numerik. Hasil keluaran berupa bitstring optimal dan energi terendah mutlak dari prosedur eksak ini dicatat secara permanen untuk memvalidasi tingkat akurasi konvergensi model hibrida kuantum.

\begin{algorithm}[ht]
\caption{Pencarian Solusi Eksak Brute Force}
\label{alg:brute_force}
\begin{algorithmic}[1]
\Procedure{BruteForceValidation}{$h, J, C, N, K$}
    \State $best\_energy \gets \infty$
    \State $best\_x \gets \text{None}$
    \State $all\_combinations \gets \text{GenerateCombinations}(N, K)$
    \For{each $x \in all\_combinations$}
        \State $E \gets \text{HitungEnergiIsing}(x, h, J, C)$ \Comment{Kalkulasi Energi Hamilton}
        \If{$E < best\_energy$}
            \State $best\_energy \gets E$
            \State $best\_x \gets x$
        \EndIf
    \EndFor
    \State \Return $best\_x, best\_energy$
\EndProcedure
\end{algorithmic}
\end{algorithm}

